\section{Portfolio Implementation}

\subsection{Portfolio Results}
Table~\ref{tab:portfolio} presents the primary portfolio outcomes under the base cost scenario (31bp round-trip).

\begin{table}[htbp]
\centering
\caption{Portfolio performance under base cost scenario (31bp round-trip).}
\label{tab:portfolio}
\begin{tabular}{lcccc}
\toprule
\textbf{Portfolio} & \textbf{Net Annual} & \textbf{Sharpe} & \textbf{Max DD} & \textbf{Excess} \\
\midrule
U100 RAW (base31) & +6.8 & +0.31 & +46.3 & +6.9 \\
U50 RESID (base31) & -0.9 & -0.04 & +61.3 & +4.9 \\
U50 EW benchmark & -5.9 & -0.26 & --- & --- \\
\bottomrule
\end{tabular}
\end{table}

The U100 portfolio achieved positive net returns but with extreme concentration in low-liquidity stocks. The U50 portfolio failed to produce positive absolute returns, despite generating positive benchmark-relative excess.

\subsection{The Excess-Return Trap}
U50 RESID generates +4.9 annual excess over a declining benchmark of -5.9. Positive excess over a poor benchmark is not evidence of investment skill.

\subsection{Cost Analysis}
Transaction costs accounted for approximately 4.2 percentage points of annual return drag. Primary drivers: stamp tax (5bp on sells, ~3.2\% p.a.) and commission plus slippage (~1.0\% p.a.).
